% ANNEXES

\chapter{Annexes}
\label{ch:annexes}

\section*{Annexe A — Manuel utilisateur (extrait)}

\subsection*{A.1 — Déposer une plainte (espace citoyen)}

\begin{enumerate}
    \item Se connecter à l'espace citoyen avec son adresse électronique. Un code à usage unique est envoyé par courriel ; il est valable dix minutes.
    \item Choisir \textbf{Déposer une plainte}, puis la \textbf{nature de l'infraction} dans la liste proposée.
    \item Renseigner le lieu des faits par la cascade \textbf{Région $\rightarrow$ Département $\rightarrow$ Arrondissement $\rightarrow$ Quartier}. Le commissariat compétent s'affiche automatiquement.
    \item Rédiger le \textbf{récit des faits}. Le bouton microphone permet de dicter la déclaration plutôt que de la saisir.
    \item Indiquer l'\textbf{identité du mis en cause} si elle est connue ; cette étape peut être passée.
    \item Répondre aux \textbf{questions complémentaires} propres à la nature déclarée. Les questions déjà couvertes par le récit ne sont pas posées.
    \item Vérifier l'\textbf{aperçu}, signer, puis valider. Le \textbf{récépissé} au format PDF, portant le numéro de dossier et un QR code de vérification, est téléchargeable immédiatement.
\end{enumerate}

\figopt{ecran_depot_nature.png}{0.58}{Étape 1 — nature de l'infraction et cascade de localisation}{fig:man_depot_nature}

\figopt{ecran_depot_apercu.png}{0.58}{Étape 5 — aperçu de la déclaration avant signature}{fig:man_depot_apercu}

\subsection*{A.2 — Suivre son dossier}

Depuis \textbf{Mes plaintes}, le plaignant consulte l'état de chacun de ses dossiers. Les statuts successifs sont : \textit{Reçu}, \textit{Audition}, \textit{En instruction}, \textit{Décision}, \textit{Transmis}, \textit{Clôturé}. Toute convocation reçue apparaît dans le dossier concerné.

\figopt{ecran_suivi.png}{0.58}{Espace citoyen — suivi de l'avancement d'un dossier}{fig:man_suivi}

\subsection*{A.3 — Coter un dossier (espace commissaire)}

Depuis le tableau de bord, sélectionner un dossier au statut \textit{Reçu}, cliquer sur \textbf{Coter}, désigner l'enquêteur destinataire et confirmer la priorité proposée. La cotation est inscrite au journal d'audit et notifiée à l'enquêteur.

\figopt{ecran_cotation.png}{0.58}{Espace commissaire — cotation d'un dossier à un enquêteur}{fig:man_cotation}

\subsection*{A.4 — Dresser un procès-verbal (espace enquêteur)}

Ouvrir un dossier coté, puis choisir \textbf{Nouveau PV} et sélectionner l'audition concernée. Une proposition rédigée à partir de la déclaration s'affiche dans l'éditeur. \textbf{La relecture est obligatoire} : corriger, compléter, puis valider et signer. Le PV n'est enregistré qu'après validation explicite.

\figopt{ecran_pv.png}{0.58}{Espace enquêteur — éditeur de procès-verbal, proposition révisable}{fig:man_pv}

\section*{Annexe B — Captures d'écran complémentaires}

\figopt{ecran_connexion.png}{0.58}{Écran de connexion — choix du rôle et authentification}{fig:cap_connexion}

\figopt{ecran_recepisse.png}{0.58}{Récépissé de dépôt de plainte généré au format PDF}{fig:cap_recepisse}

\figopt{ecran_convocation.png}{0.58}{Convocation émise par l'enquêteur}{fig:cap_convocation}

\figopt{ecran_qr.png}{0.58}{Vérification d'une pièce par son QR code}{fig:cap_qr}

\section*{Annexe C — Extraits de code}

\subsection*{C.1 — Dépôt d'une plainte et rattachement au commissariat}

\begin{lstlisting}[language=Python,caption={Endpoint Flask de dépôt d'une plainte},label={lst:depot}]
@app.route("/api/plaintes", methods=["POST"])
@jwt_required()
def deposer_plainte():
    """Enregistre une plainte et la rattache au commissariat competent."""
    donnees = PlainteSchema().load(request.get_json())
    citoyen_id = get_jwt_identity()

    # Le ressort territorial determine le commissariat : le quartier
    # des faits, et non le domicile du plaignant, fait autorite.
    commissariat = Commissariat.query.filter_by(
        quartier=donnees["quartier"]
    ).first_or_404(description="Aucun commissariat pour ce quartier")

    plainte = Plainte(
        citoyen_id=citoyen_id,
        commissariat_id=commissariat.id,
        nature=donnees["nature"],
        recit=donnees["recit"],
        statut=StatutPlainte.RECU,
        score_completude=calculer_score(donnees),
    )
    db.session.add(plainte)
    db.session.flush()          # obtient l'id avant l'ecriture du journal

    journaliser(plainte.id, "DEPOT", auteur=citoyen_id)
    db.session.commit()

    return jsonify(PlainteSchema().dump(plainte)), 201
\end{lstlisting}

\subsection*{C.2 — Rédaction assistée du procès-verbal}

\begin{lstlisting}[language=Python,caption={Appel au modèle génératif (Llama 3.1 via Groq), côté serveur},label={lst:ia}]
from groq import Groq

client = Groq(api_key=os.environ["GROQ_API_KEY"])   # jamais en dur

def proposer_proces_verbal(plainte):
    """Retourne une PROPOSITION de PV. La validation par l'enqueteur
    reste obligatoire : cette fonction n'ecrit rien en base."""
    consigne = (
        "Tu rediges un projet de proces-verbal d'audition a partir de la "
        "declaration ci-dessous. Style administratif, troisieme personne. "
        "N'invente aucun fait absent de la declaration : en cas de donnee "
        "manquante, ecris [A COMPLETER PAR L'ENQUETEUR]."
    )
    reponse = client.chat.completions.create(
        model="llama-3.1-8b-instant",
        messages=[
            {"role": "system",  "content": consigne},
            {"role": "user",    "content": f"Nature : {plainte.nature}\n"
                                           f"Declaration : {plainte.recit}"},
        ],
        temperature=0.2,        # bas : on veut de la fidelite, pas du style
    )
    return reponse.choices[0].message.content
\end{lstlisting}

\subsection*{C.3 — Cloisonnement des accès au niveau du SGBD}

\begin{lstlisting}[language=SQL,caption={Politique d'accès par ligne sur la table des plaintes},label={lst:rls}]
-- Un plaignant ne voit que ses propres dossiers.
CREATE POLICY plaintes_select_citoyen ON plaintes
  FOR SELECT USING (citoyen_id = current_user_id());

-- Un enqueteur ne voit que les dossiers qui lui ont ete cotes.
CREATE POLICY plaintes_select_enqueteur ON plaintes
  FOR SELECT USING (enqueteur_id = current_user_id());

-- Un commissaire voit tous les dossiers de son commissariat.
CREATE POLICY plaintes_select_commissaire ON plaintes
  FOR SELECT USING (commissariat_id = current_user_commissariat());
\end{lstlisting}

\section*{Annexe D — Fiches de tests}

\begin{table}[H]
    \centering
    \begin{tabularx}{\textwidth}{|l|>{\raggedright\arraybackslash}X|>{\raggedright\arraybackslash}X|l|}
        \hline
        \textbf{Réf.} & \textbf{Cas de test} & \textbf{Résultat attendu} & \textbf{Verdict} \\
        \hline
        T-01 & Dépôt d'une plainte pour vol simple, les 5 étapes renseignées & Dossier créé au statut \textit{Reçu}, récépissé PDF délivré & Conforme \\
        \hline
        T-02 & Dépôt avec identité du mis en cause laissée vide & Dépôt accepté, champ marqué « inconnu » & Conforme \\
        \hline
        T-03 & Sélection d'un quartier hors ressort couvert & Message d'erreur explicite, dépôt bloqué & Conforme \\
        \hline
        T-04 & Récit mentionnant déjà l'heure et les témoins & Questions correspondantes non posées & Conforme \\
        \hline
        T-05 & Cotation d'un dossier à un enquêteur, priorité \textit{Haute} & Dossier visible chez l'enquêteur, cotation journalisée & Conforme \\
        \hline
        T-06 & Consultation par un citoyen du dossier d'un tiers & Accès refusé & Refusé comme attendu \\
        \hline
        T-07 & Appel de l'API sans jeton d'authentification & Réponse HTTP 401 & Refusé comme attendu \\
        \hline
        T-08 & Proposition de PV puis validation sans relecture & Validation possible mais tracée au nom de l'enquêteur & Conforme, point de vigilance \\
        \hline
        T-09 & Vérification du QR code d'un récépissé & Renvoi vers le dossier correspondant & Conforme \\
        \hline
        T-10 & Déclaration dictée puis relue à l'écran & Texte transcrit, corrections manuelles nécessaires & Conforme, limite connue \\
        \hline
    \end{tabularx}
    \caption{Fiches de tests de recette}
    \label{tab:fiches_tests}
\end{table}

\section*{Annexe E — Glossaire technique}

\begin{table}[H]
    \centering
    \begin{tabularx}{\textwidth}{|l|X|}
        \hline
        \textbf{Terme} & \textbf{Définition} \\
        \hline
        API REST & Interface de programmation suivant le style architectural REST, permettant au frontend et au backend de communiquer par requêtes HTTP standardisées. \\
        \hline
        Cotation & Acte par lequel le commissaire désigne l'enquêteur chargé d'une affaire et en fixe la priorité. \\
        \hline
        Procès-verbal (PV) & Acte écrit par lequel un officier de police judiciaire consigne des déclarations ou des constatations. \\
        \hline
        Récépissé & Document remis au plaignant attestant du dépôt de sa plainte et portant le numéro de dossier. \\
        \hline
        IA générative & Système capable de produire un texte structuré à partir d'une entrée en langage naturel. Dans PlainteCam, elle ne produit que des propositions soumises à validation. \\
        \hline
        Architecture trois-tiers & Modèle séparant la présentation (frontend), la logique métier (backend) et la persistance (base de données). \\
        \hline
        Jeton JWT & Jeton signé transmis à chaque requête pour vérifier l'identité et le rôle de l'appelant, sans session côté serveur. \\
        \hline
        ORM & Technique de mapping objet-relationnel permettant de manipuler la base via des objets Python plutôt que par du SQL écrit à la main. \\
        \hline
        WSGI & Interface standard entre un serveur web et une application Python, définie par la PEP 3333. \\
        \hline
        Politique d'accès par ligne & Règle appliquée par le SGBD lui-même, restreignant les lignes visibles selon l'utilisateur connecté. \\
        \hline
        Agile Scrum & Cadre de gestion de projet itératif organisant le développement en sprints de durée fixe. \\
        \hline
    \end{tabularx}
    \caption{Glossaire des termes techniques et métier}
    \label{tab:glossaire}
\end{table}
